% ============================================================
% §2.9 · Apéndices Técnicos — Manual Técnico K-QGMIP
% ============================================================

\section{Apéndices Técnicos}

% ─────────────────────────────────────────────────────────────────────────────
\subsection{Demostraciones}
% ─────────────────────────────────────────────────────────────────────────────

\subsubsection*{Proposición 1 — Monotonicidad $\varphi(k+1) \geq \varphi(k)$}

\textbf{Enunciado.} Sea $\Pi^{(k)}$ la $k$-partición producida por E4 o C4
para un sistema dado. Entonces $\varphi(k+1) \geq \varphi(k)$ para todo
$k \geq 1$.

\textbf{Demostración.}
La función objetivo es
\[
  \Phi(\Pi) = \mathrm{EMD}\!\left(p(s_{t+1}),\;
              \bigotimes_{P_i \in \Pi} p_{P_i}(s_{t+1}^{(P_i)})\right).
\]
Ambas estrategias son \emph{divisivas}: $\Pi^{(k+1)}$ se obtiene de
$\Pi^{(k)}$ partiendo exactamente una parte $P_j$ en dos partes
$A \cup B = P_j$, $A \cap B = \emptyset$.
Sea $\mathcal{M}_\Pi$ el conjunto de distribuciones de la forma
$\bigotimes_{P_i} p_{P_i}$ para particiones más finas o iguales que $\Pi$.
Como $\Pi^{(k+1)}$ refina $\Pi^{(k)}$, se tiene
$\mathcal{M}_{\Pi^{(k+1)}} \subseteq \mathcal{M}_{\Pi^{(k)}}$.
La EMD-Effect es la distancia (norma $L_1$) entre la distribución real
$p(s_{t+1})$ y el producto tensorial.
Si el espacio de factorizaciones disponibles se achica
($\mathcal{M}_{\Pi^{(k+1)}} \subseteq \mathcal{M}_{\Pi^{(k)}}$),
la distancia mínima no puede disminuir:
\[
  \Phi\!\left(\Pi^{(k+1)}\right)
    = \mathrm{EMD}\!\left(p,\, \mathcal{M}_{\Pi^{(k+1)}}\right)
    \geq \mathrm{EMD}\!\left(p,\, \mathcal{M}_{\Pi^{(k)}}\right)
    = \Phi\!\left(\Pi^{(k)}\right). \qquad \square
\]

\subsubsection*{Proposición 2 — Regresión exacta $k=2$ de E4}

\textbf{Enunciado.} $\mathrm{KGeoMIP}(k{=}2) \equiv \mathrm{GeoMIP}$, es
decir, producen la misma bipartición con $|\Delta\varphi| < 10^{-9}$.

\textbf{Demostración.}
En la FASE 2 del pseudocódigo de E4 (ver §A.2), se llama directamente a
\texttt{GeoMIP.aplicar\_estrategia()} sobre $V$ completo y se retorna su
\texttt{Solution} sin modificación.
Para $k=2$ el bucle de la FASE 3 no ejecuta ninguna iteración
($j = 3\ldots k = 2$ es vacío).
Luego la partición retornada por E4 en $k=2$ es exactamente la partición
de GeoMIP, y $\Phi$ se calcula con la misma función
\texttt{\_calcular\_phi\_total} sobre esa misma bipartición. $\square$

\subsubsection*{Proposición 3 — Aditividad de la EMD-Effect}

\textbf{Enunciado.}
Dada una $k$-partición $\Pi = \{P_1, \ldots, P_k\}$,
la función objetivo se descompone aditivamente como
$\Phi(\Pi) = \sum_{i=1}^{k} \mathrm{costo}(P_i)$,
donde $\mathrm{costo}(P_i) = \sum_{d \in P_i} \left| p_d - \bar{x}_d(P_i) \right|$.

\textbf{Demostración.}
La EMD-Effect bajo la métrica $L_1$ es
$\mathrm{EMD}(u, v) = \sum_d |u_d - v_d|$.
La distribución reconstruida $v = \bigotimes_{P_i} p_{P_i}$ asigna a cada
nodo $d \in P_i$ su marginal $\bar{x}_d(P_i)$ (promedio sobre el estado
de los demás nodos del bloque en el estado inicial).
Entonces:
\[
  \Phi(\Pi)
    = \sum_d \left| p_d - \bar{x}_d(P_{[d]}) \right|
    = \sum_{i=1}^{k} \sum_{d \in P_i} \left| p_d - \bar{x}_d(P_i) \right|
    = \sum_{i=1}^{k} \mathrm{costo}(P_i).
\]
Esta aditividad justifica el cálculo incremental $\Delta\Phi$ :
al bipartir $P_j$ en $A$ y $B$, el incremento exacto es
$\Delta\Phi = \mathrm{costo}(A) + \mathrm{costo}(B) - \mathrm{costo}(A \cup B)
= \Phi(\Pi^{(k+1)}) - \Phi(\Pi^{(k)})$. $\square$

% ─────────────────────────────────────────────────────────────────────────────
\subsection{Detalles Algorítmicos Adicionales}
\label{sec:apendice_algoritmos}
% ─────────────────────────────────────────────────────────────────────────────

\subsubsection*{A.1 — Pseudocódigo completo de KQNodes C4 con MinHeap}

El siguiente pseudocódigo describe el criterio C4 implementado en
\texttt{KQNodes.aplicar\_estrategia}:

\begin{verbatim}
Entrada: V (nodos del subsistema), k, oracle Queyranne
Salida:  Particion Pi de k partes, Phi*

FASE 1: Inicializar
  Pi <- { V }
  (A*, B*, phi_loc) <- QNodes(V)           # Queyranne sobre V completo
  MinHeap.insertar((phi_loc, A*, B*, V))   # clave: phi_loc minima

FASE 2: Refinar k-1 veces
  Para j = 1 hasta k-1:
    (phi_sel, A_sel, B_sel, P_sel) <- MinHeap.extraer_min()
    Pi <- (Pi - {P_sel}) U {A_sel, B_sel}
    Para hijo en {A_sel, B_sel} con |hijo| >= 2:
      (A', B', phi') <- QNodes(hijo)
      MinHeap.insertar((phi', A', B', hijo))

FASE 3: Calcular Phi* exacto
  Phi* <- EMD(p(s_{t+1}), kron_{Pi_i in Pi} p_{Pi_i})
  Retornar (Pi, Phi*)
\end{verbatim}

\textbf{Complejidad}: $\mathcal{O}(k \cdot D^3 \cdot N)$ con caché del oracle.
El heap ejecuta $k{-}1$ extracciones y a lo sumo $2(k{-}1)$ inserciones.
Queyranne sobre cada parte de tamaño $D' \leq D$ cuesta $\mathcal{O}(D'^3 \cdot N)$.

\subsubsection*{A.2 — Pseudocódigo completo de KGeoMIP E4}

El siguiente pseudocódigo describe la estrategia E4 implementada en
\texttt{KGeoMIP.aplicar\_estrategia}:

\begin{verbatim}
Entrada: V (nodos), k, tabla T, estrategia_corte in {auto, exhaustivo, guiado_S}
Salida:  Particion Pi de k partes, Phi*

FASE 0: k=1 -> retornar ({V}, 0)

FASE 1: Construir S (una sola vez por subsistema)
  S[i,j] <- similitud causal entre nodos i, j via T   # O(n^2 * 2^n)

FASE 2: k=2 -> ancla GeoMIP
  (A_raiz, B_raiz, Phi_2) <- GeoMIP.find_mip(V, T)
  Si k = 2: retornar ({A_raiz, B_raiz}, Phi_2)
  # Contraste: comparar con MejorCorte(V) bajo modelo k;
  # el menor DeltaPhi gana (empate: favor ancla GeoMIP)
  Pi <- { A_raiz, B_raiz }
  MinHeap.insertar((DeltaPhi_raiz, A_raiz, B_raiz, V))

FASE 3: Para j = 3 hasta k:
  (DeltaPhi_sel, A_sel, B_sel, P_sel) <- MinHeap.extraer_min()
  Pi <- (Pi - {P_sel}) U {A_sel, B_sel}
  Para hijo en {A_sel, B_sel} con |hijo| >= 2:
    (A', B', DeltaPhi') <- MejorCorte(hijo, T, S, estrategia_corte)
    MinHeap.insertar((DeltaPhi', A', B', hijo))

FASE 4: Phi* <- EMD(p(s_{t+1}), kron_{Pi_i in Pi} p_{Pi_i})
  Retornar (Pi, Phi*)

MejorCorte(P, T, S, modo):
  Si modo = "auto":
    Si 2^(|P|-1)-1 <= 4096: modo <- "exhaustivo"
    Si no:                   modo <- "guiado_S"
  Si modo = "exhaustivo":
    Evaluar todas las biparticiones de P; retornar argmin DeltaPhi
  Si modo = "guiado_S":
    Generar hasta 20 candidatos por afinidad cruzada en S
    Evaluar DeltaPhi para cada candidato; retornar el minimo
\end{verbatim}

\subsubsection*{A.3 — Desempate determinístico en E4}

Cuando dos o más partes tienen el mismo $\Delta\Phi$ (diferencia $< 10^{-9}$),
el MinHeap aplica el siguiente orden lexicográfico:

\begin{enumerate}
  \item \textbf{Orden BFS de GeoMIP}: preferir el corte en el orden en que
        GeoMIP evaluaría sus candidatos (heurística 1 por índice de variable,
        luego heurística 2 por nivel BFS).
        Garantiza reproducir exactamente el desempate de GeoMIP en $k=2$.
  \item \textbf{Menor cardinalidad}: preferir la parte de menor $|P_i|$.
  \item \textbf{Índice léxico}: preferir la parte cuyo menor elemento tiene
        menor índice.
\end{enumerate}

La clave de prioridad en el heap es la tupla
$(\Delta\Phi,\; \text{bfs\_order},\; |P_i|,\; \min(\text{índices}(P_i)))$,
cuya comparación lexicográfica en Python implementa el desempate sin
estado externo ni semilla aleatoria.

\subsubsection*{A.4 — Memoización canónica de NCube.marginalizar }

El método \texttt{NCube.marginalizar(ejes)} computa la marginalización
de un hipercubo sumando sobre los ejes especificados.
La optimización introduce dos cambios:

\begin{enumerate}
  \item \textbf{Clave canónica por máscara de bits}:
        en lugar de \texttt{tuple(ejes)}, la clave del caché es
        \texttt{int(máscara) = ejes $\cap$ dims} expresado como bitmask.
        Esto hace que llamadas con distintos órdenes de ejes o con ejes
        fuera de \texttt{dims} apunten a la misma entrada.

  \item \textbf{Caché de instancias NCube completas}:
        el caché almacena la instancia \texttt{NCube} resultante, no solo el
        array de datos.
        Como \texttt{NCube} es inmutable (\texttt{frozen=True}),
        es seguro compartir la misma instancia entre llamadores.
        Cadenas incrementales (p.ej., el oracle de Queyranne pidiendo
        $A$, luego $A \cup \{d\}$) reutilizan los nodos ya calculados
        en vez de recomputarlos.
\end{enumerate}

\textbf{Impacto medido}: el hot path de marginalizaciones repetidas
pasa de $583\,\text{ms}$ a $13\,\text{ms}$ para $10^4$ llamadas con $n=10$
(factor $45\times$); los barridos KGeoMIP/KQNodes mejoran un 6--12\,\%.

% ─────────────────────────────────────────────────────────────────────────────
\subsection{Resultados Experimentales Complementarios}
% ─────────────────────────────────────────────────────────────────────────────

\subsubsection*{Comparación entre muestras A y B — KGeoMIP $k=3$}

La Tabla~\ref{tab:comparacion_muestras} consolida los resultados de KGeoMIP
para las dos muestras de red ejecutadas: N10A (muestra A) y N15B (muestra B).

\begin{table}[H]
\centering
\begin{tabular}{llrrrrrr}
  \toprule
  \textbf{Config.} & \textbf{Muestra} & \textbf{Casos} &
  $\varphi_{\min}$ & $\varphi_{\mathrm{avg}}$ & $\varphi_{\max}$ &
  $t_{\mathrm{avg}}$ (s) & $t_{\mathrm{total}}$ (s) \\
  \midrule
  KGeoMIP $n=10$, $k=3$ & A & 49 & 0.0283 & 1.0156 & 3.3848 & 0.0252 &  1.23 \\
  KGeoMIP $n=15$, $k=3$ & B & 50 & 0.0004 & 0.0074 & 0.0209 & 0.4556 & 22.78 \\
  \bottomrule
\end{tabular}
\caption{Comparación entre muestras A ($n=10$) y B ($n=15$) para KGeoMIP E4, $k=3$.
  Los valores de $\varphi$ en N15B son dos órdenes de magnitud menores que en N10A,
  lo que refleja las diferencias en la distribución causal de cada red.}
\label{tab:comparacion_muestras}
\end{table}

\textbf{Observación}: el caso más costoso de N15B es el subsistema completo
(caso 1: ABCDEFGHIJKLMNO $\times$ ABCDEFGHIJKLMNO, $t=2.09\,\text{s}$),
mientras que los subsistemas de 7--8 nodos terminan en $<0.1\,\text{s}$.
La estrategia \texttt{"auto"} activa \texttt{guiado\_S} para bloques de
$|P| > 13$, lo que mantiene el tiempo total bajo 23\,s para 50 casos.

\subsubsection*{Benchmark de gap de optimalidad — KGeoMIP E4}

La Tabla~\ref{tab:gap_geomip} muestra el gap de optimalidad antes y después
de la optimización ($\Delta\Phi$ incremental exacto + raíz consistente con el
modelo $k$) sobre un conjunto de casos de prueba con $n \in \{5,6,8\}$.

\begin{table}[H]
\centering
\begin{tabular}{lcrr}
  \toprule
  \textbf{Caso} & $k$ & \textbf{Gap antes} & \textbf{Gap después} \\
  \midrule
  N5A, estado \texttt{10000} & 4 & 0.375 & 0.250 \\
  N5A, estado \texttt{11111} & 4 & 0.125 & 0.125 \\
  N6A, estado \texttt{100000} & 3 & 0.406 & \textbf{0 (exacto)} \\
  N6A, estado \texttt{100000} & 4 & 0.250 & \textbf{0 (exacto)} \\
  N8A, estado \texttt{10000000} & 3 & 1.000 & \textbf{0 (exacto)} \\
  N8A, estado \texttt{10000000} & 4 & 1.000 & \textbf{0 (exacto)} \\
  \bottomrule
\end{tabular}
\caption{Gap de optimalidad $\varphi_{\text{E4}} - \varphi_{\text{óptimo}}$ antes y
  después. El gap se mide vs. búsqueda exhaustiva BruteForce-k.}
\label{tab:gap_geomip}
\end{table}


% ─────────────────────────────────────────────────────────────────────────────
\subsection{Referencias}
% ─────────────────────────────────────────────────────────────────────────────

\begin{enumerate}
  \item \textbf{Albantakis, L., Barbosa, L., Findlay, G., Signorelli, C.\,M.,
        Marshall, W., Koch, C., \& Tononi, G.} (2023).
        Integrated information theory (IIT) 4.0: Formulating the properties
        of phenomenal existence in physical terms.
        \textit{PLOS Computational Biology}, 19(10), e1011465.

  \item \textbf{Kitazono, J., Kanai, R., \& Oizumi, M.} (2018).
        Efficient algorithms for searching the minimum information partition
        in integrated information theory.
        \textit{Entropy}, 20(3), 173.

  \item \textbf{Queyranne, M.} (1998).
        Minimizing symmetric submodular functions.
        \textit{Mathematical Programming}, 82(1--2), 3--12.

  \item \textbf{Oizumi, M., Albantakis, L., \& Tononi, G.} (2014).
        From the phenomenology to the mechanisms of consciousness:
        Integrated information theory 3.0.
        \textit{PLOS Computational Biology}, 10(5), e1003588.

  \item \textbf{Knuth, D.\,E.} (2011).
        \textit{The Art of Computer Programming, Volume 4A: Combinatorial
        Algorithms, Part 1.} Addison-Wesley.
        (Algoritmo de generación de particiones de Stirling, §7.2.1.5.)

  \item \textbf{Harris, C.\,R. et al.} (2020).
        Array programming with NumPy.
        \textit{Nature}, 585, 357--362.

  \item \textbf{Virtanen, P. et al.} (2020).
        SciPy 1.0: Fundamental algorithms for scientific computing in Python.
        \textit{Nature Methods}, 17, 261--272.
\end{enumerate}
